%!TEX root = ED_Lista_Fila_Pilha

\section{Fila}
% \begin{frame}
%	Em certas aplicações, tornam-se úteis estruturas de dados constituídas por  conjuntos de elementos organizados, não pelos valores dos dados, mas em função de uma determinado critério que regulamenta a entrada e a saída dos dados na estrutura.
%	\begin{itemize}
%		\item atendimento ao cliente
%		\item processamento de tarefa
%		\item etc.
%	\end{itemize}
% \end{frame}

\begin{frame}[label=notas]
	\begin{itemize}
		\item Lista FIFO (First In, First Out)
		\item Os elementos são colocados na fila e retirados (ou processados) por ordem de chegada
			\begin{itemize}
				\item Inserção \alert<1>{apenas no final} da fila
				\item Remoção \alert<1>{apenas do início} da fila
			\end{itemize}
	\end{itemize}
	\begin{center}
		\includegraphics[width=.5\textwidth]{queue}
		%source: http://commons.wikimedia.org/wiki/File:Data_Queue.svg
	\end{center}
\end{frame}

\begin{frame}[label=notas]
	Implementação? 
	\\\pause\vfill
	Uma fila permite várias operações:
	\begin{itemize}
		\item criar uma fila vazia
		\item inserir um item novo (ao final)
		\item remover um item (do início)
		\item esvaziar a fila
		\item etc.
	\end{itemize}
\end{frame}

% \begin{frame}[fragile,label=notas]
% 	\begin{center}
% 		\begin{minipage}{.45\textwidth}
% 			\begin{lstlisting}[title=Elemento de fila]
% 				typedef struct queue_item_t {
% 				  item_t item;
% 				  struct queue_item_t *next;
% 				} queue_item_t;
% 			\end{lstlisting}			
% 		\end{minipage}
% 	\end{center}
% 	\pause
% 	\begin{center}
% 		\begin{minipage}{.45\textwidth}
% 			\begin{lstlisting}[title=Fila]
% 				typedef struct queue_t {
% 				  queue_item_t *first, *last;
% 				  /* int num_items; */
% 				} queue_t;
% 			\end{lstlisting}			
% 		\end{minipage}
% 	\end{center}
% \end{frame}
% \lstset{morekeywords={queue_item_t, queue_t}}

% \begin{frame}[fragile,label=notas]
% 	\vspace{-1\baselineskip}
% 	\begin{center}
% 		\begin{minipage}{.75\textwidth}
% 			\begin{lstlisting}[title=Criar fila]
% 				queue_t *queue_new(item_t *item) {
% 				  queue_t *fila = (queue_t *)malloc(sizeof(queue_t));
% 				  if (fila) {
% 				    if(item) enqueue(fila, *item);
% 				    else (fila->first = fila->last = NULL);
% 				  }
% 				  return fila;
% 				}
% 			\end{lstlisting}
% 		\end{minipage}
% 	\end{center}
% 	\vspace{-1.5em}\pause
% 	\begin{center}
% 		\begin{minipage}{.5\textwidth}
% 			\begin{lstlisting}[title=Esvaziar fila]
% 			                	void queue_empty(queue_t *fila) {
% 			                	  while (fila->first)
% 			                	    free(dequeue(fila));
% 			                	  fila->last = NULL;
% 			                	}
% 			\end{lstlisting}			
% 		\end{minipage}
% 	\end{center}
% \end{frame}
% \lstset{morekeywords={queue_new, queue_empty}}

% \begin{frame}[fragile,label=notas]
% 	\vspace{-1\baselineskip}
% 	\begin{center}
% 		\begin{minipage}{.75\textwidth}
% 			\begin{lstlisting}[title=Inserir]
% 				void enqueue(queue_t *fila, item_t item) {
% 				  if(!fila->last) { // fila vazia
% 				    fila->first = queue_new_item(item, NULL);
% 				    fila->last = fila->first;
% 				  } else { // ao final
% 				    fila->last->next = queue_new_item(item, NULL);
% 				    fila->last = fila->last->next;
% 				  }
% 				}
% 			\end{lstlisting}
% 		\end{minipage}
% 	\end{center}
% 	\vspace{-1.5em}\pause
% 	\begin{center}
% 		\begin{minipage}{.6\textwidth}
% 			\begin{lstlisting}[title=Remover]
% 			                	queue_item_t *dequeue(queue_t *fila) {
% 			                	  queue_item_t *aux = fila->first;
% 			                	  fila->first = aux->next;
% 			                	  return aux;
% 			                	}
% 			\end{lstlisting}			
% 		\end{minipage}
% 	\end{center}
% \end{frame}
% \lstset{morekeywords={enqueue, dequeue}}

\begin{frame}
	\framesubtitle{\subsecname\- Aplicações}
	
	As filas são utilizadas quando desejamos processar itens de acordo com a 
	ordem de chegada (o primeiro a chegar será o primeiro a ser processado.
	\\\vfill\pause
	Sistemas operacionais, por exemplo, usam filas para regular a ordem em que 
	as tarefas devem receber processamento e recursos devem ser alocados a processos.
	\subitem{fila de prioridade (inserção ordenada)}
\end{frame}

\begin{frame}
	\framesubtitle{Exercícios}
	
	\begin{itemize}
		\setlength{\itemsep}{\baselineskip}
		\item Soma de inteiros (super) longos (ex5)
		\item Controle de delimitadores em uma equação (ex6)
		\item Implementação de uma fila de prioridade descendente (ex7)
		\item Manipulação de uma sequência de caracteres (ex8)
	\end{itemize}
\end{frame}